% !TEX encoding = UTF-8 Unicode
% LaTeX transcription of:
% Bernd Stellmacher, "An Application of the Amalgam Method: The 2-Local
% Structure of N-Groups of Characteristic 2 Type",
% Journal of Algebra 190 (1997), 11-67.  Article No. JA966864.
% Transcribed from the scanned PDF (visual transcription; the embedded OCR
% text layer was not used).
\documentclass[11pt]{article}

\usepackage[T1]{fontenc}
\usepackage[utf8]{inputenc}
\usepackage{amsmath,amssymb,amsthm}
\usepackage[margin=1.15in]{geometry}
\usepackage{fancyhdr}
\usepackage[colorlinks=true,linkcolor=blue,citecolor=blue]{hyperref}

\pagestyle{fancy}
\fancyhf{}
\fancyhead[L]{\footnotesize\thepage\quad BERND STELLMACHER}
\fancyhead[R]{\footnotesize 2-LOCAL STRUCTURE OF N-GROUPS}
\renewcommand{\headrulewidth}{0pt}

% Theorem-like environments.  The paper numbers all non-main statements
% (in) a single counter per section and displays them as "(x.y)" without a
% type word; the two main theorems are displayed as "THEOREM 1" and
% "THEOREM 2".
\theoremstyle{plain}
\newtheorem{mainthm}{THEOREM}
\newtheoremstyle{stmplain}{3pt}{3pt}{\itshape}{}{\normalfont}{}{.5em}{}
\theoremstyle{stmplain}
\newtheorem{stm}{}[section]
\renewcommand{\thestm}{(\arabic{section}.\arabic{stm})}

\newcommand{\N}{\mathrm{N}}

% Section heading as printed in the paper: the bare number, centered.
\newcommand{\sect}[1]{\stepcounter{section}\par\vspace{14pt}%
  \begin{center}\textbf{#1.}\end{center}\vspace{2pt}}

% Displayed line (as in the paper's plain displays).
\newcommand{\disp}[1]{\begin{center}$\displaystyle #1$\end{center}}

\title{\bfseries An Application of the Amalgam Method:\\[2pt]
  The 2-Local Structure of $\N$-Groups of\\[2pt]
  Characteristic 2 Type}
\author{Bernd Stellmacher}
\date{}

\begin{document}

\begin{center}
{\footnotesize JOURNAL OF ALGEBRA \textbf{190}, 11--67 (1997)}\\[2pt]
{\footnotesize ARTICLE NO. JA966864}\\[6pt]
{\Large\bfseries An Application of the Amalgam Method:\\
The 2-Local Structure of $\N$-Groups of\\
Characteristic 2 Type}\\[8pt]
{\large Bernd Stellmacher}\\[4pt]
{\itshape Math.\ Seminar, Christian-Albrechts-Universit\"at zu Kiel,\\
24098 Kiel, Germany}\\[6pt]
{\itshape Communicated by Gernot Stroth}\\[4pt]
{\itshape Received December 15, 1995}\\[10pt]
{\scshape Dedicated to Professor B.\ Fischer on the occasion of his\\
60th birthday}
\end{center}

\setcounter{section}{0}

% [p. 11]
\noindent In [9] Thompson classified the finite (simple) groups $H$ satisfying:

\medskip
\noindent $(\N)$ \quad \textit{Every $p$-local subgroup of $H$ is solvable, for every prime $p$.}

\medskip
\noindent Later it turned out that his proof could be used as a pattern for the
classification of the finite simple groups in general. In the course of this
later classification Thompson's result was generalized in [5], [6], and [7] to
finite (simple) groups $H$ satisfying:

\medskip
\noindent $(\N_2)$ \quad \textit{Every 2-local subgroup of $H$ is solvable.}

\medskip
\noindent The proofs of both of the above results are subdivided as to whether
the invariant $e(H)$ is small or large. Here $e(H)$ denotes the rank of a
largest elementary abelian $p$-subgroup contained in a 2-local subgroup of
$H$, where $p$ ranges over all odd primes. It is characteristic for those
proofs that the general approach used fails if $e(H)$ is small.

\noindent In the last 10 years a new method in group theory---the amalgam
method---has been developed which seems to be well suited to deal with groups
(of characteristic 2 type) where $e(H)$ is small. In contrast to other methods
in group theory it is a completely local method. That is, it focuses on the
structure of finite (local) subgroups of a group rather than the structure of
the group itself. In fact, usually neither simplicity nor finiteness of the
group is required.

% [p. 12]
This paper can be seen as an attempt to explore the reaches of the amalgam
method in general classification problems (with arbitrary $e(H)$) and how
global properties of the group in question (finiteness, simplicity, etc.) can
be used to facilitate the amalgam method. We thought that $\N$-groups, i.e.,
finite groups satisfying $(\N_2)$, would be an interesting class of groups for
studying these aspects.

\noindent The following two theorems will be proven in this paper.

\begin{mainthm}\label{thm:1}
Let $H$ be a finite group, $S_0 \in \mathrm{Syl}_2(H)$, and
$B = C_{S_0}(\Omega_1(ZJ(S_0)))$. Suppose that
\begin{enumerate}
\item[(i)] every 2-local subgroup of $H$ containing $B$ is solvable and of
  characteristic 2 type and
\item[(ii)] there exist (at least) two maximal 2-local subgroups of $H$
  containing $S_0$.
\end{enumerate}
Then $H$ is type $L_3(2)$, $Sp_4(2)$, $G_2(2)'$, ${}^2F_4(2)'$, $M_{12}$,
$\Omega_6^+(2)$, $\Omega_6^-(3)$, or $\Omega_6^+(3)$.
\end{mainthm}

\begin{mainthm}\label{thm:2}
Suppose that $H$ is an $\N_2$-group of even order and $S_0 \in \mathrm{Syl}_2(H)$.
Then one of the following holds:
\begin{enumerate}
\item[(a)] $H$ is of type $L_3(2)$, $Sp_4(2)$, $G_2(2)'$, or ${}^2F_4(2)'$.
\item[(b)] $S_0$ is a dihedral or semidihedral group.
\item[(c)] $|S_0| = 2^5$ and there exists a maximal 2-local subgroup isomorphic to $C_2 \times S_4$ in $H$.
\item[(d)] $H$ contains a strongly embedded subgroup.
\item[(e)] There exists a 2-local subgroup $U$ of $H$ such that $O_2(U) \ne 1$.
\end{enumerate}
\end{mainthm}

\noindent In Theorem~\ref{thm:1}, ``$H$ is of type $X$'' means that $H$ contains a
pair of subgroups $P_1$ and $P_2$ such that $O_2(\langle P_1, P_2 \rangle) = 1$
and $S_0 \le P_1 \cap P_2$, and the structure of $P_1$ and $P_2$ is like that
of a pair of maximal 2-local subgroups of $X_0$, where $X \le X_0 \le
\mathrm{Aut}(X)$. The precise definition for ``$H$ is of type $X$'' is given in
Section 8 [for $X = L_3(2)$, $Sp_4(2)$, $G_2(2)'$, $\Omega_6^-(3)$, and
$\Omega_6^+(3)$], Section 9 [for $X = \Omega_6^+(2)$], and Section 10 [for
$X = M_{12}$ and ${}^2F_4(2)'$].

\noindent Some remarks about the proof of Theorem~\ref{thm:1}. The amalgam method
works with a pair of subgroups $P_1$ and $P_2$ having the following property:
\begin{enumerate}
\item[(1)] $S_0 \le P_1 \cap P_2$, $O_2(P_i) \ne 1$, $i = 1,2$, and $O_2(\langle P_1, P_2 \rangle) = 1$.
\end{enumerate}

% [p. 13]
\noindent By a nice argument of Gomi [see (4.4)], such a pair of subgroups exists
with the additional property:
\begin{enumerate}
\item[(2)] $S_0$ is contained in a unique maximal subgroup of $P_i$, $i = 1, 2$.
\end{enumerate}

\noindent Property (2) yields a particularly transparent structure for
$P_1/O_2(P_1)$; see (3.3). A refinement of Gomi's argument gives a pair of
subgroups which, apart from (1) and (2), satisfies (see (4.7)):
\begin{enumerate}
\item[(3)] Either $O^2(P_i)$ is subnormal in $C_{P_i}(\Omega_1(Z(S_0)))$ or
  $\Omega_1(Z(S_0))$ is neither normal in $P_1$ nor in $P_2$.
\end{enumerate}

\noindent In the proof of (4.7) a pushing-up theorem for $SL_2(2)$ is used. This is
a special case of Baumann's pushing-up theorem [1]. The subgroup $B$ used in
hypothesis (i) of Theorem~\ref{thm:1}---the Baumann subgroup---was first used in
that paper. However, we decided to quote [8] rather than [1] [see (2.4) and
(2.5)] because the proof in [8] is short, elementary and, more important, stays
within the framework of this paper since its proof also uses the amalgam
method.

\noindent It should be pointed out that, apart from [8] and textbook material the
reader should be familiar with, the proof has been made self-contained.

\noindent The importance of (3) for the proof can only be appreciated by reading
the proofs of Sections 7--10, but some evidence can be given here.

\noindent If $\Omega_1(Z(S_0))$ is neither normal in $P_1$ nor in $P_2$, then the
normal subgroup $Z_i = \langle \Omega_1(Z(S_0))^{P_i} \rangle$ is a noncentral
$GF(2)$ $P_i$-module with $C_{S_0}(Z_i) = O_2(P_i)$, $i = 1, 2$. Hence, the
structure of $P_i/O_2(P_i)$ can be investigated by its action on $Z_i$. That
this, together with the amalgam method, is a very effective procedure can be
seen in (8.2), where this entire case is treated.

\noindent If $\Omega_1(Z(S_0))$ is normal in $P_2$, then $P_2$ acts trivially on
$Z_2$, and the action gives no further information about $P_2/O_2(P_2)$. This
lack of information from the action is compensated by the subnormality of
$O^2(P_2)$ in $C_{P_2}(\Omega_1(Z(S_0)))$. Indeed, it is mainly here where the
global structure of $H$ facilitates the application of the amalgam method.

\noindent There is another aspect of the proof which should be mentioned. There is
no subdivision of the proof according to the size of $e(H)$. But, of course, a
priori, there is no bound on $e(P_i)$, $i = 1, 2$. The easiest way to see how
the proof treats larger values of $e(P_i)$ may be by referring to an example.

\smallskip
\noindent\emph{Let $H = (E_1 \times E_2)\langle t \rangle$, where $E_1 \cong L_3(2)$,
$E_1^t = E_2$, and $t^2 = 1$. Then $H \cong L_3(2) \wr C_2$ and $e(H) = 2$ [larger
values of $e(H)$ can be produced by substituting $\langle t \rangle$ by a larger
2-group]. Let $S_0 \in \mathrm{Syl}_2(H)$, $t \in S_0$, and $T = S_0 \cap
F^*(H)$, and let $L_1, L_2 \le E_1T$ such that $T = L_1 \cap L_2$ and
$O^2(L_i) \cong A_4$, $i = 1, 2$. Then $P_1 = \langle L_1, t \rangle$ and
$P_2 = \langle L_2, t \rangle$ are the only maximal 2-local subgroups of $H$
containing $S_0$. They both are solvable and satisfy (1)--(3).}

% [p. 14]
\noindent We now forget about the global structure of $H$ and only use the
structure of $P_1$ and $P_2$. More precisely, we only use the following
properties:
\begin{enumerate}
\item[(*)] $[\Omega_1(Z(S_0)), O^2(L_i), O^2(L_i^t)] = 1$ for $i = 1, 2$.
\item[(**)] $C_{Z(O_2(P_1))}(O^2(L_1^t)) \cap C_{Z(O_2(P_2))}(O^2(L_2^t)) \cap Z(T) \ne 1$.
\end{enumerate}

\noindent Then $\langle L_1^t, L_2^t \rangle \le C_H(x)$ for suitable
$1 \ne x \in Z(T)$, and it is fairly easy to see that $C_H(x)$ is not solvable
[if we assume the global information, we get that $E_2 \le C_H(x)$]. Hence,
$H$ is not an $\N_2$-group.

\noindent The crucial observation is that the Baumann subgroup $B$ (of $S_0$) is
contained in $L_1$ and $L_2$; in fact $B = T$ in our example. In general,
property (*) can be established if $B \le L_1 \cap L_2$ and $B$ is neither
normal in $L_1$ nor $L_2$; see (3.10). This then allows one to use the above
argument in fairly general situations. Usually the element $x$ can be found in
\disp{[\Omega_1(Z(S_0)), O^2(L_1)] \cap [\Omega_1(Z(S_0)), O^2(L_2)] \cap Z(T),}
but often the argument also works if $[\Omega_1(Z(S_0)), O^2(L_i)] = 1$ for
some $i$.

\noindent Most of the notation used in this paper is standard or will be defined
in the section where it occurs for the first time. Modules will be written
multiplicatively since they usually arise as normal subgroups or sections of
groups.

\smallskip
\noindent Let $Y$ be a group and $V$ be a $GF(2)$ $Y$-module. Then $Y$ operates
quadratically on $V$ if $[V, Y, Y] = 1$. An element $t$ of $Y$ induces a
transvection on $V$ if $|V/C_V(t)| = 2$, and $Y$ induces transvections on $V$
if $[V, Y] \ne 1$ and every element of $Y \setminus C_V(Y)$ induces a
transvection on $V$.

\sect{1}

\noindent In this section $G$ is a finite solvable group of even order and $V$ is a
finite $GF(2)$ $G$-module such that $C_G(V) = O_2(G) = 1$.

\smallskip
\noindent We will use the following

\medskip
\noindent\textit{Notation.} \quad $S \in \mathrm{Syl}_2(G)$, $W = O_2(G)$, and
$m(Y) = |V|(|C_V(Y)||Y|)^{-1}$ for $Y \le S$.

\smallskip
\disp{\mathcal{A}(V, S) = \{A \le S \mid A \text{ is elementary abelian and } m(A) \le 1\}.}
\disp{J(V, S) = \langle A \mid A \in \mathcal{A}(V, S) \rangle.}
\disp{B = C_S(C_V(J(V, S))).}
\disp{E = \langle J(V, S)^G \rangle.}

% [p. 15]
\noindent $\mathcal{A}(S)$ is the set of all subgroups $A$ of $S$ such that
$m(A) \le m(S)$, $C_S([C_W(A), N_S(A)]) = A$, and $C_S(C_V(A)) = A$.

\smallskip
\noindent $\Omega(W)$ is the set of all subgroups $D$ in $W$ such that $|D| = 3$
and $[V, D] = 4$.

\smallskip
\noindent The first lemma is well known and will be used in this paper without any
further reference.

\begin{stm}\label{st:1.1}
The following hold:
\begin{enumerate}
\item[(a)] $W = [W, S]C_W(S)$.
\item[(b)] $W = \langle C_W(a) \mid a \in S^{\#} \rangle$ if $S$ is noncyclic
  and abelian.
\item[(c)] $W = \langle C_W(S_0) \mid |S/S_0| = 2 \rangle$ if $S$ is elementary
  abelian.
\item[(d)] $S$ is elementary abelian if $[V, S, S] = 1$.
\end{enumerate}
\end{stm}

\begin{stm}\label{st:1.2}
Let $G = WS$. Suppose that $S$ is quadratic on $V$ and
$V = \langle C_V(S)^G \rangle$. Then
\disp{V = \langle C_V(S_0) \mid |S/S_0| = 2 \rangle.}
\end{stm}

\begin{proof}
We have
\disp{G = C_G(S)\langle [C_W(S_0), S] \mid |S/S_0| = 2 \rangle.}
Let $H$ be a fixed subgroup of index 2 in $S$, $N = C_W(H)$ and
$V_0 = \langle C_V(S_0) \mid |S/S_0| = 2 \rangle$. It suffices to show that
$[N, S]$ normalizes $V_0$.

\smallskip
\noindent Let $S_0$ be any subgroup of index 2 in $S$. Then the quadratic action of
$S$ on $V$ gives
\disp{[S, C_V(S_0), N] \le C_V(H) \quad\text{and}\quad [C_V(S_0), N, S] \le C_V(H).}
Hence the 3-subgroup lemma implies $[N, S, C_V(S_0)] \le C_V(H) \le V_0$.
\end{proof}

\begin{stm}\label{st:1.3}
Let $x$ be an involution of $G$ and $F = [W, x]$. Suppose that $F$ is a
$p$-group and $|V/C_V(x)| \le 4$. Then one of the following holds:
\begin{enumerate}
\item[(a)] $[V, F] = 4$ and $F \cong C_3$.
\item[(b)] $[V, F] = 2^4$, $|V/C_V(x)| = 4$, and $F \cong C_3$, $C_5$, or
  $C_3 \times C_3$.
\item[(c)] $[V, F] = 2^6$, $|V/C_V(x)| = 4$, $[Z(F), x] = 1$, and $F$ is extra
  special of order $3^3$.
\end{enumerate}
\end{stm}

\begin{proof}
We proceed by induction on $|F||V|$. Then $G = F\langle x \rangle$ and $V =
[V, F]$; in particular $C_V(F) = 1$. Note that $|V/C_V(x)| = |[V, x]|$.

\smallskip
\noindent Assume first that there exists $1 \ne a \in Z(F)$ such that $a^x = a^{-1}$
and $[V, a] = 2^4$. Then $[V, x] \le [V, a]$, and $F = [F, x]$ implies $V =
[V, a]$. Hence, $F$ is a subgroup of $GL_4(2)$ ($\cong A_8$), and (a) or (b)
follows.

% [p. 16]
\noindent Assume next that there exists $1 \ne a \in Z(F)$ such that $a^x = a^{-1}$
and $[V, a] = 4$. Note that $C_F(V/[V, a]) = \langle a \rangle$ since $C_G(V) =
1$. Set $\bar{V} = V/[V, a]$ and $\bar{G} = G/\langle a \rangle$. By induction
either $F = \langle a \rangle$ or $\bar{F} \cong C_3$ and $F = Z(F)$ since
$|\bar{V}/C_{\bar{V}}(\bar{x})| \le 2$. The first case gives (a). In the second
case there exists $b \in Z(F)$ such that $b^x = b^{-1}$ and $[V, b] = 2^4$, and
the above argument applies.

\smallskip
\noindent We may assume now that $[Z(F), x] = 1$. Let $z \in Z(F)$. If $[V, x, z] =
1$, then $[V, F, x] \le C_V(z)$ and $[V, x, F] \le C_V(z)$. Hence also
\disp{[F, x, V] = [F, V] = V \le C_V(z)}
and $z = 1$. This shows that $[V, x] = 4$, $Z(F) \cong C_3$, and $V = [V,
Z(F)]$. In particular, $V$ can be regarded as a $GF(4)$ $G$-module by the
action of $Z(F)$.

\smallskip
\noindent Assume next that $F = \langle a, b \rangle$, where $a^x = a^{-1}$ and $b^x
= b^{-1}$. Then $V = [V, a][V, b]$ and $|V| \le 2^6$ since $[V, a] = [V, b] =
2^4$ and $[V, x] \le [V, a] \cap [V, b]$. Hence, $G$ is a subgroup of $SL_3(4)$
and (c) follows.

\smallskip
\noindent Assume finally that there exist $a, b \in F$ so that $a^x = a^{-1}$, $b^x
= b^{-1}$, $[a, b] \ne 1$, and $F \ne \langle a, b \rangle$. Set $F_0 =
\langle a, b \rangle$ and $V_0 = [V, F_0]$. By induction $|V_0| = 2^6$ and
$F_0$ is extra special of order $3^3$. Note that the structure of $SL_3(4)$
gives $V_0 \ne V$ and $C_F(Z(F_0)) \ne F$. Let $F_1 = C_F(Z(F_0))$ and $F_2 =
N_F(F_1)$. Again by induction $F_1 = C_{F_1}(x)F_0$ and $[Z(F_1), x] = 1$. It
follows that $[F_2, x] \ne F_0$ and so $[F_2, x, Z(F_1)] \ne 1$. On the other
hand, $[x, Z(F_1), F_2] = [Z(F_1), F_2, x] = 1$ which contradicts the
3-subgroup lemma.
\end{proof}

\begin{stm}\label{st:1.4}
Let $\Omega(W) = \{F_1, \ldots, F_r\}$ and $W_0 = \langle F_1, \ldots, F_r \rangle
\le O_3(G)$. Then
\begin{enumerate}
\item[(a)] $W_0 = F_1 \times \cdots \times F_r$ and
\item[(b)] $V = V_0 \times V_1 \times \cdots \times V_r$, where $V_0 = C_V(W_0)$
  and $V_i = [V, F_i]$.
\end{enumerate}
\end{stm}

\begin{proof}
Note that $[V, \langle F_1, F_2 \rangle] \le 2^4$. Hence, the structure of
$GL_4(2)$ gives (a) and (b).
\end{proof}

\begin{stm}\label{st:1.5}
Let $U$ be a subgroup of $S$. Then
\begin{enumerate}
\item[(a)] $|C_V(U)/C_V(S)| = m(S)m(U)^{-1}|S/U|$.
\end{enumerate}
Moreover, if $S$ is elementary abelian, then
\begin{enumerate}
\item[(b)] $S = \langle A \mid A \in \mathcal{A}(S), |A| = 2 \rangle$,
\item[(c)] $W = \langle C_W(A) \mid A \in \mathcal{A}(S), |S/A| = 2 \rangle$,
\item[(d)] $W = \langle C_W(A) \mid A \in \mathcal{A}(S), |A| = 2 \rangle$ if
  $|S| \ge 4$,
\item[(e)] $m(S) \ge 1$.
\end{enumerate}
\end{stm}

% [p. 17]
\begin{proof}
Claim (a) is obvious. Let $S$ be elementary abelian and $A$ be any subgroup of
index 2 in $S$. From (a) we get
\disp{|C_V(A)/C_V(S)| = 2m(S)m(A)^{-1}.}
Hence either $C_V(A) = C_V(S)$ or $m(A) \le m(S)$.

\smallskip
\noindent Let $Y = [C_W(A), S]$. Assume that $C_V(A) = C_V(S)$. Then $Y$ centralizes
$C_V(S)$ and $V = C_V(Y) \times [V, Y]$. Since $[V, Y]$ is $S$-invariant and
$C_V(S) \le C_V(Y)$, we get that $[V, Y] = 1$ and $C_W(A) = C_W(S)$. Note
further that $C_S(Y) \ne A$ also implies $C_W(A) = C_W(S)$. Hence (c) follows.

\smallskip
\noindent From (c) we get that $S = \langle A \mid A \in \mathcal{A}(S), |S/A| = 2
\rangle$. If $\mathcal{A}(A) \subseteq \mathcal{A}(S)$ for every $A \in
\mathcal{A}(S)$ with $|S/A| = 2$, then (b) and (d) follow by induction on
$|S|$. Hence, to prove (b) and (d), it suffices to prove the above inclusion
$\mathcal{A}(A) \subseteq \mathcal{A}(S)$.

\smallskip
\noindent Let $A \in \mathcal{A}(S)$ with $|S/A| = 2$ and $A_0 \in
\mathcal{A}(A)$. Set $Y_0 = [C_W(A_0), S]$ and $V_0 = C_V(A_0)$. Since $Y \le
Y_0$ and $C_V(A) \le V_0$ we get that $C_S(Y_0) \le C_S(Y) = A$ and $C_S(V_0)
\le C_S(C_V(A)) = A$; i.e., $C_S(Y_0) = C_A(Y_0)$ and $C_S(V_0) = C_A(V_0)$.
Now $A_0 \in \mathcal{A}(A)$ implies that $A_0 \in \mathcal{A}(S)$.

\smallskip
\noindent To prove (e), let $A \in \mathcal{A}(S)$ and $|A| = 2$. Then (a) implies
that $m(A) \ge 1$ and thus $m(S) \ge 1$.
\end{proof}

\begin{stm}\label{st:1.6}
Let $V^* = [V, W]$. Suppose that $S$ is elementary abelian, $W = [W, S]$, and
$m(S) \le m(Y)$ for every subgroup $Y \ne 1$ of $S$. Then one of the following
holds:
\begin{enumerate}
\item[(a)] $|S| = 2$ and $m(S) > 1$,
\item[(b)] $m(S) = 2$, $G \cong SL_2(2) \times SL_2(2)$, $|V^*| = 2^4$, and
  $[V^*, S, S] \ne 1$,
\item[(c)] $W = F_1 \times \cdots \times F_r$, where $F_i \in \Omega(W)$,
  $[V, S, S] = 1$, and $|V^*/C_{V^*}(S)| = 2|S|$,
\item[(d)] $WS = E_1 \times \cdots \times E_r$, where $E_i \cong SL_2(2)$ and
  $E_i' \in \Omega(W)$, $[V, S, S] = 1$, and $|V^*/C_{V^*}(S)| = |S|$.
\end{enumerate}
\end{stm}

\begin{proof}
If $|S| = 2$, then (a) resp.\ (d) follows easily. Hence, we may assume that
$|S| \ge 4$.

\smallskip
\noindent Let $\mathcal{A}_2(S) = \{A \in \mathcal{A}(S) \mid |A| = 2\}$ and $A \in
\mathcal{A}_2(S)$, and let $V_A = C_V(A)$ and $W_A = [C_W(A), S]$. Note that
$m(A) = m(S)$ and that $C_{W_A}(V_A) = 1$ by the $P \times Q$ lemma. From
(1.5)(a) we get
\disp{|V_A/C_{V_A}(S)| = |S/A| = |S/C_S(V_A)|.}
Hence, by induction on $|S|$ we may assume that $W_AS/A = \bar{E}_1 \times
\cdots \times \bar{E}_s$ and $V_A = U_0 \times U_1 \times \cdots \times U_s$,
where $\bar{E}_i \cong SL_2(2)$, $U_0 = C_{V_A}(W_AS)$, $U_i = [V_A, \bar{E}_i]$,
and $|U_i| = 4$. In particular, $W_A$ is a 3-group. Since $W$ is solvable and
$[W, S] = W$, (1.5)(d) implies that $W$ is a 3-group.

% [p. 18]
\noindent Suppose that there exists $E_i$ such that $[V, E_i'] \nleq V_A$. Let
$A_i = S \cap E_i$, $U = [V, E_i]$, and $S_0 = C_S(U)$. Then $[U, A_i, A_i] \ne
1$ and $|A_i| = 4$. Moreover, $|V_A/C_{V_A}(A_i)| = 2$ and so for $a \in A_i^{\#}$
either $C_V(a) \ge V_A$ or $|C_V(a)/C_{V_A}(a)| = 2$ since $m(\langle a \rangle)
\ge m(S) = m(A)$. Since $U \nleq V_A$ we get
\disp{|C_V(a)/C_{V_A}(a)| = |[V, a]/[V_A, a]| = 2 \qquad \text{for } a \in A_i \setminus A.}
It follows that $|U| = 2^4$ and $[U, a] = 4$ for $a \in S \setminus S_0$. In
particular, $[U, a, s] = [U, a, A]$ for $a \in A_i \setminus A$ and $s \in
C_S(E_i) \setminus S_0$. This gives $[U, s] = [U, A]$ by the action of $E_i$,
and $C_S(E_i) = A$. Hence $S = A_iS_0$.

\smallskip
\noindent Assume that $S_0 \ne 1$. Then by (1.5)(a),
\disp{|C_V(S_0)/C_V(S)| = 4m(S)m(S_0)^{-1},}
while $|U/C_V(S)| = 8$, a contradiction. Hence $S_0 = 1$ and $|S| = 4$.

\smallskip
\noindent Let $W_1 = C_W(U)$. Then $W_1 \le C_W(S)$ since $[V, S] \le U$. As shown
above $[C_W(X), S] \cong C_3$ for $X \in \mathcal{A}_2(S)$. Hence $[C_W(X), S,
W_1] = 1$ and $W_1 \le Z(W)$ by (1.5)(d). It follows that $[V, W_1, W_1] = 1$
since $[V, W_1, S] = 1$ and $W = [W, S]$. This implies that $W_1 = 1$ and $W
\cong C_3 \times C_3$. Now (b) is easy to verify.

\smallskip
\noindent Suppose now that for every $A \in \mathcal{A}_2(S)$, $[V, E_i'] \le V_A$
for $i = 1, \ldots, s$. Together with (1.4) and (1.5)(d) this shows that $W =
F_1 \times \cdots \times F_r$, $F_i \in \Omega(W)$, and $S \le N_G(F_i)$. Hence
$[V^*, S, S] = 1$. We now choose $A \in \mathcal{A}_2(S)$ with the additional
property the $s$ is maximal; i.e., $|[W, A]|$ is minimal.

\smallskip
\noindent As above, for $i = 1, \ldots, s$ there exists $A \le A_i \le S$ such that
$|A_i| = 4$, $|V_A/C_{V_A}(A_i)| = 2$, and $[W_A, A_i] = F_i$. Let $i$ be fixed
and $A = \langle a \rangle$. Among all subgroups of order 2 in $A$, we choose
$B_i$ such that $B_i \ne A$ and $[W, B_i]$ is minimal with that property. Let
$B_i = \langle x \rangle$. Since $a$ inverts $[W, A]$ the minimality of $[W,
x]$ gives $C_{[W, A]}(x) = 1$. Hence, there exists $F_k \in \Omega(W)$ such that
$[F_k, x] = 1$ and $[F_k, A] = F_k$. It follows that $[V, F_k] \le C_V(x) \nleq
V_A$ and $m(B_i) \le m(A)$; in particular $B_i \in \mathcal{A}_2(S)$. Thus, the
minimality of $[W, A]$ gives $[W, A] = [W, B_i] = 3$ or $9$. Let $B_0 = A$ and
$V_0 = C_V(W)$, and note that $S = \langle B_0, \ldots, B_s \rangle$.

\smallskip
\noindent Assume that $[W, A] = 3$. Then $WS \cong SL_2(2) \times \cdots \times
SL_2(2)$, $|S| = 2^s$, $|V^*/C_{V^*}(S)| = |S|$, and $|V^*/C_{V^*}(B_i)| = 2$.
It follows that
\disp{|V/C_V(S)| = |S||V_0/C_{V_0}(S)| = |S|m(S)}
and
\disp{|V/C_V(B_i)| = 2|V_0/C_{V_0}(B_i)| = 2m(B_i).}
Hence $C_{V_0}(S) = C_{V_0}(B_i)$ and $[V, S, S] = 1$ since $m(B_i) = m(S)$,
and (d) follows.

% [p. 19]
\noindent Assume now that $[W, A] = 9$. Then $|S| = 2^{r-1}$ since $W_AS/A \cong
SL_2(2) \times \cdots \times SL_2(2)$ and $|V^*/C_{V^*}(B_i)| = 4$. Now as above
$C_{V_0}(S) = C_{V_0}(B_i)$ and $[V, S, S] = 1$, and (c) follows.
\end{proof}

\begin{stm}\label{st:1.7}
Suppose that $J(V, S) \ne 1$. Then the following hold:
\begin{enumerate}
\item[(a)] $W = E' \times C_W(E)$.
\item[(b)] $B = J(V, S)$ and $J(V, S) \in \mathcal{A}(V, S)$.
\item[(c)] $E = E_1 \times \cdots \times E_n$ and $V = V_0 \times V_1 \times
  \cdots \times V_n$, where $V_0 = C_V(E)$, $E_i \cong SL_2(2)$, $V_i = [V,
  E_i]$, and $|V_i| = 4$ for $i = 1, \ldots, n$.
\end{enumerate}
\end{stm}

\begin{proof}
Let $A \in \mathcal{A}(V, S)$. Then $m(A) \le 1$, and (1.5)(e) gives $m(A) = 1$.
Now (1.6)(d) and again (1.5)(e) imply that $[W, A], A = E_1 \times \cdots \times
E_s$ and $V = C_V([W, A], A) \times V_1 \times \cdots \times V_s$, where $E_i
\cong SL_2(2)$, $V_i = [V, E_i]$, and $|V_i| = 4$. Since $E_i = [W, E_i \cap
A](E_i \cap A)$ each of the subgroups $E_i$ is normal in $W_A$.

\smallskip
\noindent Let $\Omega^*$ be the set of all subgroups $F$ of $G$ such that $F \cong
SL_2(2)$, $[V, F] = 4$, and $F$ is normal in $WF$, and let $E_0 = \langle F \mid
F \in \Omega^* \rangle$. Then $E_0 = E_1 \times \cdots \times E_r$, where
$\Omega^* = \{E_1, \ldots, E_r\}$. Note that $E_0$ is normal in $G$.

\smallskip
\noindent Let $S_0 = S \cap E_0$. As we have seen above $J(V, S) \le S_0$. On the
other hand, $|V/C_V(S_0)| \le |S_0|$ and thus $S_0 \in \mathcal{A}(V, S)$; in
particular $S_0 = J(V, S)$ and $E_0 = E$.

\smallskip
\noindent It remains to prove $B = S_0$. Clearly, $B$ normalizes every $E_i \in
\Omega^*$. Hence $B = S_0C_B(E)$. However, $C_B(E)$ centralizes $V$ and thus
$B = S_0$.
\end{proof}

\sect{2}

\noindent In this section $G$ is a finite solvable group of even order and
$C_G(O_2(G)) \le O_2(G)$.

\smallskip
\noindent\textit{Notation.} \quad $S \in \mathrm{Syl}_2(G)$, $Z = \Omega_1(Z(S))$, $V =
\langle Z^G \rangle$, and $\bar{G} = G/C_G(V)$. $\mathcal{A}(S)$ is the set of
all elementary abelian subgroups of maximal order of $S$, $J(S) = \langle A
\mid A \in \mathcal{A}(S) \rangle$, $\bar{Z} = \Omega_1(Z(J(S)))$, $B =
C_S(\bar{Z})$, and $L = \langle B^G \rangle$.

\begin{stm}\label{st:2.1}
$O_2(\bar{G}) = 1$.
\end{stm}

\begin{proof}
This follows directly from the definition of $V$.
\end{proof}

\begin{stm}\label{st:2.2}
Let $\bar{E} = [O_2(\bar{G}), \bar{J}(\bar{S})]\bar{J}(\bar{S})$. Suppose that
$\bar{J}(\bar{S}) \ne 1$. Then the following hold:
\begin{enumerate}
\item[(a)] $\bar{E}$ is a normal subgroup of $\bar{G}$.
\item[(b)] $\bar{J}(\bar{S}) = \bar{B}$.
\item[(c)] $\bar{E} = \bar{E}_1 \times \cdots \times \bar{E}_n$, where $\bar{E}_i
  \cong SL_2(2)$.
\item[(d)] $V = V_0 \times \cdots \times V_n$, where $V_0 = C_V(\bar{E})$, $V_i =
  [V, E_i]$, and $|V_i| = 4$ for $i = 1, \ldots, n$.
\end{enumerate}
\end{stm}

\begin{proof}
Let $A \in \mathcal{A}(S)$. Then the maximality of $|A|$ yields
\disp{|VC_A(V)| = |V||C_A(V)||V \cap A|^{-1} \le |A|.}
It follows that $|V/C_V(A)| \le |\bar{A}|$ and $\bar{A} \in \mathcal{A}(V,
\bar{S})$ (for definition see Section 1); in particular $\bar{J}(\bar{S}) \le
J(V, \bar{S})$. Now (1.6)(d) and (1.7) applied to $\bar{G}$ and $V$ give (c),
(d), and, together with an easy Frattini argument, also (a).

\smallskip
\noindent Note that $\bar{B}$ normalizes $\bar{E}_i$ for $i = 1, \ldots, n$ and
centralizes $V_0$ since $V_0 \le \bar{Z}$. It follows that
$C_{\bar{B}}(\bar{E}') = C_{\bar{B}}(V) = 1$. Hence $\bar{J}(\bar{S}) \le
\bar{B}$ implies that $\bar{J}(\bar{S}) = \bar{B}$.
\end{proof}

\begin{stm}\label{st:2.3}
Suppose that $O_2(G) = C_S(V)$. Then $B \in \mathrm{Syl}_2(L)$.
\end{stm}

\begin{proof}
Assume first that $[V, J(S)] = 1$. Then $V \le Z$ and thus $B \le C_S(V) =
O_2(G)$. Hence, $B$ is normal in $G$.

\smallskip
\noindent Assume now that $[V, J(S)] \ne 1$. Let $E = \langle J(S)^G \rangle$. We
apply (2.2). Then $\bar{J}(\bar{S}) = \bar{B}$ and $L = E\bar{B}$. Hence, it
suffices to show that $O_2(E) \le B$. Let $Z_0 = \Omega_1(Z(J(O_2(E))))$ and $A
\in \mathcal{A}(S)$. Note that by (1.5)(e), $C_A(V)V \in \mathcal{A}(S)$.
Hence $\bar{Z} \le Z_0$ and $Z_0 \in C_{Z_0}(A)V$. It follows that $[Z_0, E]
\le V$; in particular, $\bar{Z}V$ is normal in $E$. Now the structure of
$\bar{E}$ and its operation on $\bar{V}$ shows that there exists $x \in E$ so
that $E = \langle J(S), J(S)^x \rangle C_E(V)$ and $\bar{Z}V = Z^*V$, where
$Z^* = \bar{Z} \cap \bar{Z}^x$ and $\bar{Z} = Z^*(\bar{Z} \cap V)$. Moreover,
$C_E(V)$ centralizes $\bar{Z}V$ and so $C_E(V) = C_E(Z^*V)O_2(E)$. Since
$O_2(E)$ normalizes $Z^*$, we conclude that $Z^*$ is normal in $E$. However,
now $[Z^*, J(S)] = 1$ implies $[Z^*, E] = 1$, and $O_2(E) \le C_S(\bar{Z}) = B$.
\end{proof}

\begin{stm}\label{st:2.4}
Suppose that the following hold:
\begin{enumerate}
\item[(i)] No nontrivial characteristic subgroup of $S$ is normal in $G$.
\item[(ii)] $S$ is contained in a unique maximal subgroup of $G$.
\end{enumerate}
Then $[O_2(G), O^2(G)] \le V$.
\end{stm}

\begin{proof}
By (i) neither $Z$ nor $J(S)$ is normal in $G$; i.e., $[V, J(S)] \ne 1$.
Moreover, (ii) implies that $C_S(V) = O_2(G)$. We apply (2.2) and (2.3).

\smallskip
\noindent Let $B \le L_1 \le L$ so that $\bar{L}_1/O_2(\bar{L}_1) \cong SL_2(2)$ and
$|L_1|$ is minimal with that property. Then $B \in \mathrm{Syl}_2(L_1)$ and
$L_1/O_2(L_1) \cong D_{2 \cdot 3}$. Moreover, by (ii), $\langle L_1, S \rangle =
G$ and thus no nontrivial characteristic subgroup of $B$ is normal in $L_1$.
Hence, [8] implies that $[O_2(L_1), O^2(L_1)] \le V$. Let $T$ be a Hall
$2'$-subgroup of $C_L(O_2(L)/V)$. Then $L_1 \le TS$ and $TS = G$. It follows
that $[O_2(G), O^2(G)] \le V$.
\end{proof}

% [p. 21]
\begin{stm}\label{st:2.5}
Suppose that $G$ satisfies the hypothesis of (2.4) and $\bar{G} \cong SL_2(2)$.
Let $T$ be a subgroup of $\mathrm{Aut}(S)$ of odd order. Then $\langle V^{\tau}
\mid \tau \in T \rangle$ is a normal subgroup of $G$ in $O_2(G)$.
\end{stm}

\begin{proof}
This is [8, 3.5].
\end{proof}

\sect{3}

\noindent In this section $G$ is a finite group of even order, $S$ is a nontrivial
2-subgroup of $G$, and $U$ is a subgroup of $G$ containing $S$.

\smallskip
\noindent\textit{Notation.} \quad $\mathcal{L}(U, S) = \{E \le U \mid S \in
\mathrm{Syl}_2(E),\, O_2(E) \ne 1, \text{ and } S \ne O_2(E)\}$;
$\mathcal{P}(U, S) = \{E \in \mathcal{L}(U, S) \mid S$ is contained in a unique
maximal subgroup of $E\}$; $\mathcal{P}^*(U, S) = \{E \in \mathcal{P}(U, S)
\mid O^2(E)$ is subnormal in $L$ for some maximal element $L \in
\mathcal{L}(U, S)\}$. $\Phi_2(U)$ is the inverse image of $\Phi(U/O_2(U))$ in
$U$. If $U = G$, we also write $\mathcal{L}(S)$, $\mathcal{P}(S)$, and
$\mathcal{P}^*(S)$ instead of $\mathcal{L}(G, S)$, $\mathcal{P}(G, S)$, and
$\mathcal{P}^*(G, S)$, respectively.

\begin{stm}\label{st:3.1}
Let $S \in \mathrm{Syl}_2(G)$ and $\mathcal{F}$ be the set of all subgroups $U$
of $G$ such that $S \le U$ and
\begin{enumerate}
\item[(*)] $S \ne O_2(U)$ and $S$ is contained in a unique maximal subgroup of
  $U$.
\end{enumerate}
Then either $\mathcal{F} = \emptyset$ and $S = O_2(G)$, or $O^2(G) = \langle P
\mid P \in \mathcal{F} \rangle$.
\end{stm}

\begin{proof}
Let $G$ be a minimal counterexample and $G_0 = \langle P \mid P \in
\mathcal{F} \rangle$, and let $M$ be any maximal subgroup of $G$ containing
$S$. Then either $M = N_G(S)$ or, by induction, $M = N_M(S)(M \cap G_0)$.
Since $N_G(S)$ normalizes $G_0$ we conclude that $G_0N_G(S)$ is the unique
maximal subgroup of $G$ containing $S$, but now $G \in \mathcal{F}$. Since $G =
O^2(G)N_G(S) = O^2(G)M$ it follows that $G = O^2(G)$, and $G$ is not a
counterexample.
\end{proof}

\begin{stm}\label{st:3.2}
Let $L \in \mathcal{L}(S)$. Then $O^2(L) = \langle P \mid P \in
\mathcal{P}(L, S) \rangle$.
\end{stm}

\begin{proof}
This is a direct consequence of (3.1).
\end{proof}

\begin{stm}\label{st:3.3}
Let $P \in \mathcal{P}(S)$ and $B$ be the maximal subgroup of $P$ containing
$S$, and let $P_0$ be the largest normal subgroup of $P$ contained in $B$.
Suppose that $P$ is solvable. Then
\begin{enumerate}
\item[(a)] $O^2(P/O_2(P))$ is a $p$-group for some odd prime $p$,
\item[(b)] $O^2(P)/P_0$ is an irreducible $S$-module,
\item[(c)] $P_0/O_2(P) = \Phi(O^2(P/O_2(P)))$.
\end{enumerate}
\end{stm}

\begin{proof}
Let $\bar{P} = P/O_2(P)$. The existence of Hall subgroups in $\bar{P}$ shows
that $\bar{P}$ is a $(2, p)$-group, where $p$ is some odd prime, and the
Frattini argument shows that $P_0$ is a $p$-group.
\end{proof}

% [p. 22]
\noindent Clearly, $\Phi(O^2(\bar{P})) \le \bar{P}_0$, and Maschke's theorem gives
$\bar{P}_0 = \Phi(O^2(\bar{P}))$ and (b).

\begin{stm}\label{st:3.4}
Let $P \in \mathcal{P}(S)$ and $T$ be a normal subgroup of $S$. Suppose that $P$
is solvable. Then either $T \le O_2(P)$ or $[O^2(P), T] = O^2(P)$.
\end{stm}

\begin{proof}
Since $[O^2(P), T]$ is normal in $P$, the claim follows from (3.3).
\end{proof}

\begin{stm}\label{st:3.5}
Let $P \in \mathcal{P}(S)$ and $N$ be a normal subgroup of $P$ in $O_2(P)$.
Suppose that $P$ is solvable and $[N, O_2(P) \cap O^2(P)] = 1$. Then either
$[Z(S), O^2(P)] \ne 1$ or $[N, O^2(P)] = 1$.
\end{stm}

\begin{proof}
Note that $[N, O^2(P)] \le Z(O_2(O^2(P)))$ and that $[N, O^2(P), O^2(P)] = 1$
implies $[N, O^2(P)] = 1$. Hence we may assume that $N \le O^2(P)$ and $N$ is
abelian. By Maschke's theorem $N = C_N(O^2(P)) \times [N, O^2(P)]$. Since
$[N, O^2(P)]$ is $S$-invariant, the claim follows.
\end{proof}

\begin{stm}\label{st:3.6}
Let $P \in \mathcal{P}(S)$, $\bar{P} = P/O_2(P)$, and $Z$ and $T$ be two normal
subgroups of $S$, and let $A$ be a subgroup of $S$ satisfying $\Phi(A) \le
O_2(P)$ and $a \in A \setminus O_2(P)$. Suppose that $P$ is solvable, $Z \le
O_2(P)$, and $T \nleq O_2(P)$. Then there exists $x \in P$ so that for $L =
\langle A, A^x \rangle$ the following hold:
\begin{enumerate}
\item[(a)] $\bar{L} = \bar{E} \times \bar{A}_0$, where $\bar{E} \cong D_{2p^n}$
  and $A = \langle a \rangle A_0$.
\item[(b)] $O^2(L) \nleq \Phi_2(O^2(P))$.
\item[(c)] Any two elements in $Z^L$ are interchanged by an involution of $L$.
\item[(d)] $O^2(L) \le [O^2(L), T]$.
\end{enumerate}
\end{stm}

\begin{proof}
Let $F = [O^2(P), a]$. If $F \le \Phi_2(O^2(P))$, then $a \in O_2(P)$ which is
not the case. Hence $F \nleq \Phi_2(O^2(P))$. Let $F_0 \le F$ such that
\begin{enumerate}
\item[(*)] $F_0$ is $A$-invariant and $F_0 \nleq \Phi_2(O^2(P))$.
\end{enumerate}
Among all such $F_0$ satisfying (*) we choose $F_0$ such that first
$|F_0[F_0, T]|$ and then $|F_0|$ is minimal. Since $F = \langle C_F(\bar{A}_0)
\mid |\bar{A}/\bar{A}_0| = 2 \rangle$ there exists $A_0 \le A$ such that
$|A/A_0| = 2$ and $F_0 \le C_F(A_0)$ by the minimal choice of $F_0$. Note that
$\bar{F}_0 \nleq \Phi(\bar{F})$ and so $\bar{a} \notin \bar{A}_0$.

\smallskip
\noindent Choose $\bar{z} \in \bar{F}_0 \setminus \Phi(O^2(\bar{P}))$ such that
$\bar{z}^{\bar{a}} = \bar{z}^{-1}$. The minimality of $\bar{F}_0$ gives
$\bar{F}_0 = \langle \bar{z} \rangle$. Set $L = \langle A, A^x \rangle$. Then
(a) and (b) hold. Since $O_2(L)$ normalizes every element in $Z^L$, the
structure of $L/O_2(L)$ also gives (c).

\smallskip
\noindent By (3.4), $\bar{F}_0 \le [\bar{F}_0, \bar{T}]\Phi(O^2(\bar{P}))$. Hence
$F_1 = [F_0, T, a] \nleq \Phi_2(O^2(P))$. Now the minimality of $F_0[F_0, T]$
gives $F_0 \le [F_0, T]$, and (d) follows.
\end{proof}

\begin{stm}\label{st:3.7}
Let $L \in \mathcal{L}(S)$ and $\bar{L} = L/O_2(L)$. Suppose that $L$ is
solvable. Then
\disp{[F(\bar{L}), \bar{S}] = \langle O^2(\bar{P}) \mid P \in \mathcal{P}^*(L, S) \rangle.}
\end{stm}

% [p. 23]
\begin{proof}
Let $F$ be the inverse image of $[F(\bar{L}), \bar{S}]$ in $L$. By (3.2), $FS =
\langle P \mid P \in \mathcal{P}(FS, S) \rangle$. Moreover, $O^2(P)$ is
subnormal in $L$ for every $P \in \mathcal{P}(FS, S)$; i.e., $P \in
\mathcal{P}^*(L, S)$.

\smallskip
\noindent Now let $P \in \mathcal{P}^*(L, S)$. Then $O^2(P)$ is subnormal in $L$
since $L$ is the unique maximal element of $\mathcal{L}(L, S)$. It follows that
$O^2(\bar{P}) \le F(\bar{L})$ and so $P \in \mathcal{P}(FS, S)$ and
$\mathcal{P}^*(L, S) = \mathcal{P}(FS, S)$.
\end{proof}

\begin{stm}\label{st:3.8}
Let $P_1, P_2 \in \mathcal{P}(S)$ and $H = \langle P_1, P_2 \rangle$, and let $N$
be a normal subgroup of $H$ which is maximal (with respect to inclusion) such
that $O^2(P_i) \nleq N$ for $i = 1, 2$. Suppose that $P_1$ and $P_2$ are
solvable. Then there exists a normal series $Q \le N \le H_0 \le H_1 \le H$
such that
\begin{enumerate}
\item[(a)] $Q = S \cap N$ and $Q$ is the largest subgroup of $S$ which is normal
  in $H$,
\item[(b)] $H_0/N$ is a minimal normal subgroup of $H/N$ and $O^2(P_j) \le H_0$
  for some $j \in \{1, 2\}$,
\item[(c)] $H_1 = H_0$ or $H_1 = H_0O^2(P_k)$, where $O^2(P_k) \nleq H_0$,
\item[(d)] $H = SH_1$.
\end{enumerate}
\end{stm}

\begin{proof}
Let $Q^* = S \cap O_2(H)$. Clearly, $Q^*$ is the largest normal subgroup of $H$
in $S$. Hence $Q^* \le N$ by the maximality of $N$; i.e., $Q^* \le S \cap N =
Q$. On the other hand, by (3.4), $Q = O_2(P_i) \cap N$ since $O^2(P_i) \nleq
N$. Hence $Q = Q^*$.

\smallskip
\noindent Let $H_0$ be the inverse image of a minimal normal subgroup of $H/N$. Then
there exists $j \in \{1, 2\}$ such that $O^2(P_j) \le H_0$. Let $\{1, 2\} =
\{j, k\}$. If $O^2(P_k) \le H_0$, then $H_0S = H$. If $O^2(P_k) \nleq H_0$,
then $H_0O^2(P_k)S = H$.
\end{proof}

\begin{stm}\label{st:3.9}
Let $P_1, P_2$ and $H$ be as in (3.8), and let $S \le T \in \mathrm{Syl}_2(H)$
and $Q = S \cap O_2(H)$. Suppose that $H$ is solvable, $\Omega_1(Z(S)) \le Q$,
and $J(S) = J(T)$. Then one of the following holds:
\begin{enumerate}
\item[(a)] $[\Omega_1(Z(S)), O^2(P_i), O^2(P_j)] = 1$ for $i \ne j$, or
\item[(b)] $C_S(\Omega_1(Z(J(S)))) \le O_2(P_i)$ for some $i \in \{1, 2\}$, or
\item[(c)] $[\Omega_1(Z(S)), O^2(P_1)] = [\Omega_1(Z(S)), O^2(P_2)]$.
\end{enumerate}
\end{stm}

\begin{proof}
Let $Z = \Omega_1(Z(S))$, $V = \langle Z^H \rangle$, $\bar{H} = H/C_H(V)$, and
$B = C_S(\Omega_1(Z(J(S))))$. We may assume that $B \nleq O_2(P_i)$ and
$O^2(P_i) \nleq C_H(V)$ for $i = 1, 2$. We apply (3.8) with $C_H(V) \le N$.

\smallskip
\noindent Assume that $J(S) \le C_H(V)$. Then $J(S) \le Q$ and $V \le Z(J(S))$.
Hence $B \le C_H(V)$ and so $B \le Q \le O_2(P_i)$, $i = 1, 2$, a contradiction.

\smallskip
\noindent Assume now that $J(S) \nleq C_H(V)$. Then there exists $i \in \{1, 2\}$
such that $J(S) \nleq O_2(P_i)$; i.e., $O^2(P_i) \le \langle J(S)^{P_i} \rangle$
by (3.4). Let $\bar{F} = O_2(F(\bar{H})) \cdot O^2(\bar{P}_i)J(\bar{S})$.

% [p. 24]
\noindent Since $J(S) = J(T)$ and $\bar{S} \cap O_2(\bar{H}) = 1$ we get
$[J(\bar{S}), O_2(\bar{H})] = 1$ and thus $[\bar{F}, O_2(\bar{H})] = 1$. It
follows that
\disp{O_2(\bar{F}) \le C_{\bar{S}}(F(\bar{H})) = \bar{S} \cap O_2(\bar{H}) = 1.}
Hence, (1.7) applied to $\bar{F}$ and $\bar{V}$ gives $O^2(\bar{P}_i) \le
O_3(\bar{H})$; in particular, $\bar{H} = O_3(\bar{H})\bar{P}_j$, where $\{1, 2\}
= \{i, j\}$, and $\bar{S} \in \mathrm{Syl}_2(\bar{H})$.

\smallskip
\noindent Assume that $J(S) \nleq O_2(P_j)$. Then by the same argument, $O^2(\bar{P}_i)
\le O_3(\bar{H})$ and $\bar{H} = O_3(\bar{H})\bar{S}$. If $\bar{P}_1 \ne
\bar{P}_2$, then again (1.7) gives (a). In the other case, (c) follows.

\smallskip
\noindent Assume that $J(S) \le O_2(P_j)$. Since $B \nleq O_2(P_j)$ we get that
$[\Omega_1(Z(J(S))), O^2(P_i)] = 1$. However, now (1.7) applied to $\bar{H}$
and $\bar{V}$ shows that $[V, O^2(P_1), O^2(P_2)] = 1$, and again (a) holds.
\end{proof}

\sect{4}

\noindent In this section $G$ is a finite group of even order, $S \in
\mathrm{Syl}_2(G)$, and the following hold:
\begin{enumerate}
\item[(a)] Every 2-local subgroup of $G$ containing $S$ is solvable and of
  characteristic 2 type.
\item[(b)] $S$ is contained in two different maximal 2-local subgroups of $G$.
\end{enumerate}

\smallskip
\noindent\textit{Notation.} \quad $Z = \Omega_1(Z(S))$, $C = C_G(Z)$, $D =
\bigcap_{P \in \mathcal{P}(S)} O_2(P)$, $M = N_G(D)$, $B = C_S(\Omega_1(Z(J(S))))$,

\medskip
\disp{\mathcal{P}_0(S) = \mathcal{P}^*(C, S) \cup (\mathcal{P}^*(S) \setminus \mathcal{P}(C, S)),}
\disp{\mathcal{P}_1(S) = \mathcal{P}^*(M, S) \cup (\mathcal{P}^*(S) \setminus \mathcal{P}(M, S)),}
\disp{\Lambda = \{(P_1, P_2) \mid P_i \in \mathcal{P}(S),\, O_2(\langle P_1, P_2 \rangle) = 1\}.}

\smallskip
\noindent\textit{Remark.} A subgroup similar to $D$ was used in [3].

\begin{stm}\label{st:4.1}
The following hold:
\begin{enumerate}
\item[(a)] $\mathcal{L}(S) \ne \emptyset \ne \mathcal{P}(S)$.
\item[(b)] $O_2(M) \ne 1$.
\item[(c)] $N_G(S) \le M$.
\end{enumerate}
\end{stm}

\begin{proof}
$N_G(S)$ is not the only 2-local subgroup of $G$ containing $S$. Hence (3.2)
implies (a).

\smallskip
\noindent Claim (b) holds since $Z \le D \le O_2(M)$, and (c) holds since $N_G(S)$
operates on the elements of $\mathcal{P}(S)$ by conjugation.
\end{proof}

% [p. 25]
\begin{stm}\label{st:4.2}
Let $L \in \mathcal{L}(S)$, $D_L = \bigcap_{P \in \mathcal{P}(L,S)} O_2(P)$,
and $D_L^* = \bigcap_{P \in \mathcal{P}^*(L,S)} O_2(P)$. Then $D_L = D_L^* =
O_2(L)$ and $D = \bigcap_{L \in \mathcal{M}\mathcal{L}(S)} O_2(L)$, where
$\mathcal{M}\mathcal{L}(S)$ is the set of maximal elements of
$\mathcal{L}(S)$.
\end{stm}

\begin{proof}
Let $F_0$ be the inverse image of $F(L/O_2(L))$ in $L$ and $F = [F_0, S]$. By
(3.7), $FS = \langle P \mid P \in \mathcal{P}^*(L, S) \rangle$. It follows that
$[F_0, D_L^*] \le F$ and $[F, D_L^*] \le O_2(L)$. We conclude that
\disp{O_2(L) \le D_L \le D_L^* \le O_2(L).}
and the claim follows.
\end{proof}

\begin{stm}\label{st:4.3}
$D = \bigcap_{P \in \mathcal{P}_i(S)} O_2(P)$ for $i = 0, 1$.
\end{stm}

\begin{proof}
For $\emptyset \ne \mathcal{X} \subseteq \mathcal{L}(S)$ define $O_2(\mathcal{X})
= \bigcap_{P \in \mathcal{X}} O_2(P)$. Let $M_0 = C$ and $M_1 = M$. By (4.2) it
suffices to show that $O_2\mathcal{P}_i(S) \le O_2(L)$ for every maximal
element $L \in \mathcal{L}(S)$.

\smallskip
\noindent We have $\mathcal{P}^*(L, S) \subseteq \mathcal{P}(M_i, S) \cup
(\mathcal{P}^*(S) \setminus \mathcal{P}(M_i, S))$, and by (4.2),
\disp{O_2(\mathcal{P}^*(M_i, S)) = O_2(\mathcal{P}(M_i, S)) = O_2(M_i).}
Hence
\disp{O_2(\mathcal{P}_i(S)) = O_2(\mathcal{P}(M_i, S)) \cap
O_2(\mathcal{P}^*(S) \setminus \mathcal{P}(M_i, S))}
\disp{\quad\quad\quad\quad \le O_2(\mathcal{P}^*(L, S)) = O_2(L).}
\end{proof}

\begin{stm}\label{st:4.4}
Let $P \in \mathcal{P}(S) \setminus \mathcal{P}(M, S)$ and $k \in \{0, 1\}$. Then
there exists $P^* \in \mathcal{P}_k(S)$ such that $(P, P^*) \in \Lambda$.
\end{stm}

\begin{proof}
Let $\mathcal{P}_k(S) = \{P_1, \ldots, P_n\}$ and $H_i = \langle P, P_i \rangle$,
$i = 1, \ldots, n$. We may assume that $O_2(H_i) \ne 1$ for $i = 1, \ldots,
n$. Hence by (4.2) and (4.3),
\disp{D \le \bigcap_{i=1}^n O_2(H_i) \le \bigcap_{i=1}^n O_2(P_i) = D}
and $D = \bigcap_{i=1}^n O_2(H_i)$. However, now $P \le M$, a contradiction.
\end{proof}

\begin{stm}\label{st:4.5}
Let $\mathcal{P}(S) = \mathcal{P}(M, S) \cup \mathcal{P}(C, S)$ and $P \in
\mathcal{P}(S) \setminus \mathcal{P}(M, S)$. Then there exists $P^* \in
\mathcal{P}^*(M, S)$ so that $(P, P^*) \in \Lambda$. If in addition
$\mathcal{P}^*(C, S) \subseteq \mathcal{P}(M, S)$, then $O_2(M) \le O_2(C)$
and $O_2(C) \in \mathrm{Syl}_2(O^2(P^*)O_2(C))$.
\end{stm}

\begin{proof}
By (4.4) there exists $P^* \in \mathcal{P}_1(S)$ such that $(P, P^*) \in
\Lambda$. Since $P \le C$, we have $P^* \nleq C$ and thus $P^* \in
\mathcal{P}^*(M, S)$.

\smallskip
\noindent Suppose that $\mathcal{P}^*(C, S) \subseteq \mathcal{P}(M, S)$. Then
(4.2) implies that $O_2(M) \le O_2(C)$. Since $O_2(O^2(P^*)) \le O_2(M)$, we
conclude that $O_2(C) \in \mathrm{Syl}_2(O^2(P^*)O_2(C))$.
\end{proof}

% [p. 26]
\begin{stm}\label{st:4.6}
Suppose that $\mathcal{P}(S) = \mathcal{P}(C, S) \cup \mathcal{P}(M, S)$. Then
$\mathcal{P}^*(C, S) \nsubseteq \mathcal{P}(M, S)$.
\end{stm}

\begin{proof}
Assume that $\mathcal{P}^*(C, S) \subseteq \mathcal{P}(M, S)$. By (4.5) there
exist $P \in \mathcal{P}(C, S)$ and $P^* \in \mathcal{P}^*(M, S)$ such that
$(P, P^*) \in \Lambda$ and $O_2(C) \in \mathrm{Syl}_2(O^2(P^*)O_2(C))$.

\smallskip
\noindent Let $B_0 = C_{O_2(C)}(\Omega_1(Z(J(O_2(C)))))$, $L = \langle B_0^{P^*} \rangle$,
$V = \langle Z^{P^*} \rangle$, and $\bar{L} = L/C_L(V)$. Note that $B_0$ is
normal in $P$ and thus not normal in $P^*$. Hence by (3.4), $O^2(P^*) \le L$,
and by (3.3), $O_2(L) \in \mathrm{Syl}_2(C_L(V))$ since $P^* \nleq C$. Thus
(2.3) gives $B_0 \in \mathrm{Syl}_2(L)$. Moreover, since $(P, P^*) \in \Lambda$
there is no nontrivial characteristic subgroup of $B_0$ which is normal in $L$.
Hence $LS = P^*$ and (2.2) imply
\begin{enumerate}
\item[(i)] $\bar{L} = \bar{E}_1 \times \cdots \times \bar{E}_n$,
\item[(ii)] $V = V_0 \times \cdots \times V_n$,
\end{enumerate}
where the notation is as in (2.2). Since $V_0 \cap Z$ is normalized by $P$ and
$P^*$, we get $V_0 \cap Z = 1$ and thus $V_0 = 1$.

\smallskip
\noindent Let $T$ be a Sylow 3-subgroup of $P^*$. Then by (2.4), $[O_2(P^*), T] = V$
and $O_2(P^*) = VC_{O_2(P^*)}(T)$. The Frattini argument shows that
$C_{O_2(P^*)}(T)$ is normal in $S$, and $V_0 = 1$ gives $C_{O_2(P^*)}(T) = 1$
and $O_2(P^*) = V_1 \times \cdots \times V_n$; in particular, $O_2(P^*) =
O_2(M)$.

\smallskip
\noindent Let $R_i = [V_i, B_0]$ and $\Omega = \{R_i \mid i = 1, \ldots, n\}$. Note
that $\Omega = \{R_i^s \mid s \in S\}$ since $P^* \in \mathcal{P}(S)$. Choose
$\alpha \in \mathrm{Aut}(B_0)$. If $[V_i, \alpha, V_j] \ne 1$, then $R_j =
[V_i, \alpha, V_j] = [V_i, \alpha, B_0] = R_i\alpha$. If $V_i\alpha \le V_j$,
then there exists $E_j$ so that $[V_i, \alpha, E_j \cap B_0] = R_j = R_i\alpha$.
We conclude that $\mathrm{Aut}(B_0)$ operates on $\Omega$. Since $C \le
N_G(B_0)$ we get that $C = SC_0$, where $C_0 = N_C(R_1)$.

\smallskip
\noindent Let $U$ be a Hall $2'$-subgroup of $C_0$ and $u \in U$. By (2.5),
$[V_1, V_1^u] = 1$, and $[V_i, V_1^u] \le R_1 \cap R_i = 1$ for $2 \le i$. It
follows that $\langle V_1^U \rangle \le O_2(P^*)$. Hence also
\disp{\langle \langle V_1^U \rangle^S \rangle = \langle V_1^{US} \rangle \le O_2(P^*).}
However, $US = C$ and $\langle V_1^S \rangle = O_2(P^*)$. Now $O_2(P^*) = O_2(M)$
implies that $C \le M$, and $M$ is the unique maximal 2-local subgroup
containing $S$. a contradiction.
\end{proof}

\begin{stm}\label{st:4.7}
There exists $(P, P^*) \in \Lambda$ such that $P \nleq C$ and $P^* \in
\mathcal{P}_0(S)$.
\end{stm}

\begin{proof}
Assume that $\mathcal{P}(S) \ne \mathcal{P}(C, S) \cup \mathcal{P}(M, S)$. Then
(4.4) yields the assertion. Assume that $\mathcal{P}(S) = \mathcal{P}(C, S)
\cup \mathcal{P}(M, S)$. Then (4.6) implies that $\mathcal{P}^*(C, S)
\nsubseteq \mathcal{P}(M, S)$. Now (4.7) follows from (4.5) for $P^* \in
\mathcal{P}^*(C, S) \setminus \mathcal{P}(M, S)$.
\end{proof}

% [p. 27]
\sect{5}

\noindent In this section we assume

\medskip
\noindent\textsc{Hypothesis 1.} \quad $H$ is a finite group of even order, $S_0 \in
\mathrm{Syl}_2(H)$, and the following hold:
\begin{enumerate}
\item[(i)] Every 2-local subgroup of $H$ containing $S_0$ is solvable and of
  characteristic 2 type.
\item[(ii)] $O_2(H) = 1$ and $H$ does not contain a strongly embedded subgroup.
\end{enumerate}

\begin{stm}\label{st:5.1}
There exists a nontrivial subgroup $S$ of $S_0$ and $P_1, P_2 \in \mathcal{P}(S)$
such that $O_2(\langle P_1, P_2 \rangle) = 1$ and one of the following holds:
\begin{enumerate}
\item[(a)] $S = S_0$ and $\Omega_1(Z(S))$ is neither normal in $P_1$ nor in
  $P_2$,
\item[(b)] $S = S_0$ and $P_2 \in \mathcal{P}^*(C_H(\Omega_1(Z(S))), S)$,
\item[(c)] $S \ne S_0$, $S_0$ is contained in a unique maximal 2-local subgroup
  $M$ of $H$, and
\begin{enumerate}
\item[(c$_1$)] $\Omega_1(Z(S))$ and $J(S)$ are neither normal in $P_1$ nor in
  $P_2$,
\item[(c$_2$)] $P_i \nleq M$ and $S \in \mathrm{Syl}_2(N_H(O_2(P_i)))$ for $i =
  1, 2$, and
\item[(c$_3$)] if $J(S) \le T \le S_0$, $P_1^*, P_2^* \in \mathcal{P}(T)$, and
  $O_2(\langle P_1^*, P_2^* \rangle) \ne 1$, then either $\langle P_1^*, P_2^*
  \rangle \le M$ or $J(S) = J(S_1)$ for $T \le S_1 \in \mathrm{Syl}_2(\langle
  P_1^*, P_2^* \rangle)$.
\end{enumerate}
\end{enumerate}
\end{stm}

\begin{proof}
Assume first that $S_0$ is contained in two different maximal 2-local subgroups
of $H$. Then (a) or (b) follows from (4.7).

\smallskip
\noindent Assume now that there exists a unique maximal 2-local subgroup $M$ of $H$
containing $S_0$. Since $M$ is not strongly embedded in $H$ there exists
$1 \ne Q \le S_0$ such that $N_H(Q) \nleq M$. Among all 2-local subgroups which
are not in $M$ we choose $N$ such that for $S \in \mathrm{Syl}_2(N \cap M)$
consecutively
\begin{enumerate}
\item[(i)] $|J(S)|$ is maximal,
\item[(ii)] $|S|$ is maximal.
\end{enumerate}
After conjugation in $M$ we may assume that $S \le S_0$. Since $S \ne S_0$ we
have that
\begin{enumerate}
\item[(*)] $N_H(J(S)) \le M$ and $C_H(\Omega_1(Z(S))) \le M$;
\end{enumerate}
in particular $S \in \mathrm{Syl}_2(N)$. Hence by (3.2), $\mathcal{P}(N, S)
\subsetneq \mathcal{P}(M, S)$.

\smallskip
\noindent Let $P \in \mathcal{P}(N, S) \setminus \mathcal{P}(M, S)$ and let $x \in
N_{S_0}(S) \setminus S$ with $x^2 \in S$. Set $P_1 = P$, $P_2 = P_1^x$, and
$H_0 = \langle P_1, P_2 \rangle$. Since $x \in N_H(O_2(H_0))$ and $H_0
\nsubseteq M$ the maximality of $S$ implies that $O_2(H_0) = 1$; and similarly
$S \in \mathrm{Syl}_2(N_H(O_2(P_i)))$. Together with (*), (c$_1$) and (c$_2$)
follow.

\smallskip
\noindent Now let $P_1^*$ and $P_2^*$ be as in (c$_3$) and let $H^* = \langle P_1^*,
P_2^* \rangle$ and $T \le S_1 \in \mathrm{Syl}_2(H^*)$. Assume that $J(S_1) \ne
J(S)$. Then the maximality of $J(S)$ gives $H^* \le M$. Hence (c$_3$) holds.
\end{proof}

% [p. 28]
\begin{stm}\label{st:5.2}
Let $Z_0 = \Omega_1(Z(S_0))$, $B_0 = C_{S_0}(\Omega_1(Z(J(S_0))))$, $C =
C_H(Z_0)$, and $P \in \mathcal{P}^*(C, S_0)$, and let $K$ be a subgroup of
$P$. Suppose that the following hold:
\begin{enumerate}
\item[(i)] Every 2-local subgroup of $H$ containing $B_0$ is solvable and of
  characteristic 2 type, and
\item[(ii)] $K = [K, B_0]$.
\end{enumerate}
Then $K$ is subnormal in every 2-local subgroup of $H$ containing $B_0K$.
\end{stm}

\begin{proof}
Let $D$ be a 2-subgroup of $H$ which is normalized by $KB_0$. Note that $B_0$ is
normal in $DB_0$ and by (ii), $K = O^2(K)$. Hence $[D, B_0] \le D \cap B_0$ and
\disp{[D, KB_0] = [D, \langle B_0^K \rangle] \le \langle (D \cap B_0)^K \rangle \le KB_0.}
It follows that $D \in N_H(KB_0)$, and $K = O^2(KB_0)$ yields $D \le N_H(K)$. We
have shown:
\begin{enumerate}
\item[(1)] Let $D$ be a 2-subgroup of $H$ and $KB_0 \le N_H(D)$. Then $D \le
  N_H(K)$.
\end{enumerate}
From (1) we get that $K$ is normal in $KO_2(P)$ and from (3.3) that $KO_2(P)$
is subnormal in $P$. Hence
\begin{enumerate}
\item[(2)] $K$ is subnormal in $O^2(P)$.
\end{enumerate}

\smallskip
\noindent Assume now that $K$ is a counterexample such that $|K|$ is maximal. Set
$N = N_H(K)$. By (1) there exists $T \in \mathrm{Syl}_2(N)$ such that
$O_2(C)B_0 \le T$. Let $g \in H$ such that $T \le S_0^g$. Then $g \in
N_H(B_0)$ and $Z_0^g \le Z(O_2(C))$. On the other hand, the subnormality of
$O^2(P)$ in $C$ gives $O_2(O^2(P)) \le O_2(C)$ and by (3.5), $[\Omega_1(Z(O_2(C))),
O^2(P)] = 1$. Hence $O^2(P) \le C^g$.

\smallskip
\noindent If $K \ne O^2(P)$, then the maximality of $K$ implies that $O^2(P)$ and
thus by (2) also $K$ is subnormal in $C^g$. If $K = O^2(P)$, then $S_0 \in
\mathrm{Syl}_2(N)$ and by (2), $K$ is subnormal in $C$. Hence, we have shown:
\begin{enumerate}
\item[(3)] There exists $h \in N_H(B_0)$ such that $S_0^h \cap N \in
  \mathrm{Syl}_2(N)$ and $K$ is subnormal in $C^h$.
\end{enumerate}

\smallskip
\noindent We set $T_0 = S_0^h$ and $\hat{Z}_0 = Z_0^h$, where $h$ and $S_0$ are as
in (3). Since $K$ is a counterexample there exists a subgroup $M$ such that the
following hold:
\begin{enumerate}
\item[(4)] $C_H(O_2(M)) \le O_2(M)$ and $KB_0 \le M$.
\item[(5)] $K$ is not subnormal in $M$.
\item[(6)] $K$ is subnormal in every proper subgroup of $M$ containing
  $KB_0O_2(M)$.
\end{enumerate}

\smallskip
\noindent By (1), $O_2(M) \le N$ and thus by (3), $O_2(M)B_0 \le T_0^g$ for some $g
\in N$. Then $B_0^g = B_0$ and we may assume that $T_0 = T_0^g$. Now (4) implies
that $\hat{Z}_0 \le Z(O_2(M))$.

% [p. 29]
\noindent Let $V = \langle \hat{Z}_0^M \rangle$ and $C_0 = C_M(V)$. If $K \le C_0$,
then by (3), $K$ is subnormal in $C_0$ since $C_0 \le C^h$, and thus $K$ is
subnormal in $M$ which contradicts (5).

\smallskip
\noindent Suppose that $O_2(K) \le C_0$. Then $[V, K] \le Z(O_2(K))$ by (1) and thus
$[Z(O_2(K)), K] \ne 1$ since $K \nleq C_0$. This contradicts (3.5). We have
shown:
\begin{enumerate}
\item[(7)] $O_2(K) \nleq C_0$.
\end{enumerate}

\smallskip
\noindent Let $\bar{M} = M/C_0$ and $\bar{W} = O_2(\bar{M})$. Then as above in the
proof of (1), $[\bar{W}, \bar{B}_0] \le \bar{W} \cap \bar{B}_0$ and $\bar{W} \le
N_{\bar{M}}(\bar{K})$, i.e., $W \le N_M(KC_0)$. On the other hand, $K$ is
subnormal in $KC_0$ and thus $O_2(K) \le O_2(KC_0)$. Hence
\disp{O_2(K) \cap W \le O_2(W) \le O_2(M) \le C_0.}
Since $K/O_2(K)$ is a $p$-subgroup we get that $[\bar{W}, \bar{K}] = O_2(\bar{K})
\cap \bar{W} = 1$. In particular,
\disp{[O_2(F(\bar{M})), O_2(\bar{K})] \ne 1}
by (7). Now (6) implies that $\bar{M} = O_2(F(\bar{M}))K\bar{B}_0$ and $\bar{W} \le
\bar{T}_0 \cap \bar{M}$. The definition of $\bar{V}$ implies that $\bar{W} =
1$. Since by (7), $V \nleq Z(B_0)$ we get $J(\bar{T}_0) \nleq C_0$. Hence
(2.2) gives $\bar{K} \le O_3(\bar{M})$ and $O_2(\bar{K}) = 1$, a contradiction
to (7).
\end{proof}

\smallskip
\noindent We now have set the stage for the amalgam method, which will deal with
a triple $(P_1, P_2, S)$ as in (5.1). In case (5.1)(c), $P_1$ and $P_2$ need
not be solvable or of characteristic 2 type. To get these properties we will
make a further hypothesis which will be used in most of the following sections.

\medskip
\noindent\textsc{Hypothesis 2.} \quad Hypothesis 1 holds, $S$, $P_1$, and $P_2$ are as in
(5.1), and $B = C_S(\Omega_1(Z(J(S))))$. In addition:
\begin{enumerate}
\item[(iii)] Every 2-local subgroup of $H$ containing $B$ is solvable and of
  characteristic 2 type.
\end{enumerate}

\begin{stm}\label{st:5.3}
Assume Hypothesis 2. Then $P_1$ and $P_2$ are solvable and of characteristic 2
type.
\end{stm}

\begin{proof}
Let $N = N_H(O_2(P_i))$. By (5.1), $S \in \mathrm{Syl}_2(N)$, and by Hypothesis
2, $N$ is solvable and of characteristic 2 type.
\end{proof}

\begin{stm}\label{st:5.4}
Assume Hypothesis 2. Let $B \le T \le S$, $F_1, F_2 \in \mathcal{P}(T)$, and
$H_0 = \langle F_1, F_2 \rangle$, and let $M$ be a maximal 2-local subgroup of
$H$ containing $S_0$. Suppose that $O_2(H_0) \ne 1$ and either $S = S_0$ or
$H_0 \nleq M$. Then $\Omega_1(Z(S)) \le O_2(H_0)$.
\end{stm}

\begin{proof}
Let $T_0 = T \cap O_2(H_0)$, $N = N_H(T_0)$, and $T \le S_1 \in
\mathrm{Syl}_2(H_0)$. If $S = S_0$, then obviously $J(S) = J(S_1)$; and if $S
\ne S_0$, then (5.1)(c) shows
\end{proof}

%% CONTINUE %%
\end{document}
